\documentclass[../main.tex]{subfiles}
\begin{document}
%==========================================TIÊU ĐỀ TẬP========================================
\begin{table}[h]
  \begin{adjustwidth}{-1cm}{}
    \begin{tabular}{>{\centering\arraybackslash}p{6cm}|>{\raggedright\arraybackslash}p{12.5cm}}
      \multirow{5}{6cm}
      % Cột bên trái: ÉPISODE và số tập
      {\\[-40pt]
        \flaregothic\fontsize{20pt}{20pt}\selectfont \centering
        \color{eptype}ÉPISODE\color{black}\\
        \burbank\fontsize{72pt}{72pt}\selectfont
        \color{epnum}67\color{black}\\[6pt]
        \cafeteria\fontsize{20pt}{20pt}\selectfont\color{epnum}(5-11)\color{black}
        \\[18pt]
        \flaregothic\fontsize{18pt}{18pt}\selectfont
        \color{horror} HORREUR\\
        \color{epattr} ÉPISODE PERDU\\
        \color{black}
      }
       &
      % Dòng thứ nhất: tên tập tiếng Nhật
      {\fotskip\fontsize{18pt}{18pt}\selectfont \ruby{夜}{よる}の\ruby{画}{えいが}\ruby{鑑}{かん}\ruby{賞}{しょう}\ruby{会}{かい}！
      }                                      \\[6pt]
      % Dòng thứ hai: tên tập tiếng Anh
       & {\cafeteria\fontsize{18pt}{18pt}\selectfont Overtime!}                                                                                      \\[4pt]
      % Dòng thứ ba: tên tập tiếng Việt
       & {\vnmsans\fontsize{18pt}{18pt}\selectfont Đêm xem phim kinh dị}                                                                             \\[8pt]
      % Dòng thứ tư: Nhân vật xuất hiện
       & {\caviar{\fontsize{14pt}{14pt}\selectfont Nhân vật: \emp{kuon1} \emp{achan} \emp{sora} \emp{roboco} \emp{shion} \emp{pekora} \emp{flare1}}}
      \\[8pt]
      % Dòng cuối cùng: Ngày phát hành
       & {\continuum\fontsize{16pt}{16pt}\selectfont Ngày phát hành: 09/06/2019}                                                                     \\[18pt]
    \end{tabular}
  \end{adjustwidth}
\end{table}
%=========================================TRUYỆN CHÍNH========================================
\color{black}

% Hộp tóm tắt
\txtbx[2]{{\color{vol} \uline{Tóm tắt:}} Một buổi tối mùa hè, Pekora, Roboco và Flare kéo nhau đến văn phòng để xem phim kinh dị. A-chan vẫn cắm mặt vào công việc muộn. Sora và Kuon xuất hiện, nhưng hóa ra đó là những bản sao do Shion tạo ra để quay một bộ phim giả tưởng. Khi bí mật được tiết lộ, cả nhóm cùng nhau thưởng thức một buổi marathon phim đầy tiếng cười – và Pekora lại một lần nữa sợ phim kinh dị. Đây là sản phẩm hợp tác giữa Coco, Towa, Okayu và Hokuro, với sự tham gia diễn xuất của Sora.}

Dưới bầu trời đêm giữa tháng Tám, văn phòng \hll\ chỉ còn được bao bọc bởi một màu đen tối, chỉ có duy nhất ánh đèn từ màn hình máy tính lẻ loi nơi A-chan đang chăm chú làm việc. Bên ngoài cửa sổ, những vì sao lấp lánh như những con mắt tò mò nhìn vào căn phòng đầy ánh sáng nhân tạo. Đôi tay cô thoăn thoắt lướt trên bàn phím, gõ từng dòng lệnh một cách thuần thục, nhưng trong lòng thì chỉ muốn thở dài. Cô đã làm việc suốt từ chiều đến giờ, và những con số trên màn hình bắt đầu nhòe đi trước mắt cô.

``\textit{Một ngày mới, lại thêm một ca làm đêm nữa\dots}'' cô lẩm bẩm với vẻ cam chịu, rồi đưa tay chỉnh lại cặp kính cho ngay ngắn. Cô tự hỏi tại sao mình lại đồng ý làm thêm giờ, nhưng câu trả lời luôn là giống nhau: vì có quá nhiều việc phải làm, và không đủ thời gian trong ngày.

\img{5-11a}

Cánh cửa văn phòng bật mở ra, ba bóng người bước vào. Roboco đi trước, nhẹ nhàng như thể đã quen thuộc với cảnh tượng này. Cô đã thấy A-chan thức khuya làm việc quá nhiều lần để còn ngạc nhiên: ``Em biết ngay là chị vẫn đang cắm mặt vào công việc mà,'' nàng Người máy nheo mắt cười, giọng nói có chút tinh nghịch nhưng vẫn đầy sự quan tâm, ``chị có biết bây giờ là mấy giờ rồi không? Mười một giờ đêm rồi đấy.''

Phía sau cô, Flare bước vào với dáng vẻ thư thái, đôi mắt ánh lên nét vui vẻ. Cô đang cầm một túi bim bim và một chai nước ngọt trên tay: ``Muốn xem phim thì TV ở văn phòng là số một, nên bọn em mới kéo nhau đến đây. Hy vọng chị không phiền nhé,'' cô mỉm cười, đôi tai yêu tinh khẽ lay động theo từng lời nói. ``Bọn em định tổ chức một buổi xem phim nhỏ, nhưng mà\dots''

Trong khi đó, Pekora chậm chạp bước vào sau cùng, tay ôm chặt chiếc gối nhỏ mang theo. Đôi mắt cam đảo quanh căn phòng, dáng vẻ không mấy yên tâm. Cô nhìn vào bóng tối ngoài cửa sổ, rồi nhìn vào những góc tối của căn phòng, như thể đang chờ đợi một điều gì đó nhảy ra: ``Không phải là phim kinh dị chứ? Chắc chắn không phải chứ, peko?'' giọng cô nàng lộ rõ vẻ sợ sệt, gần như van nài. ``Nếu như thế thì tôi sẽ gặp ác mộng luôn á!''

Flare liếc nhìn nàng Thỏ rồi bật cười, nhẹ nhàng vỗ vai cô như để trấn an: ``Bình tĩnh nào. Chỉ là bộ phim về một người có ngoại hình giống với một người khác đi chọc ghẹo mọi người thôi. Hoàn toàn không có ma, không có giết người, không có gì đáng sợ cả.''

Nhưng thay vì an tâm, đôi tai dài của Pekora bỗng dựng thẳng lên, nét mặt càng tái đi. Cô nhìn Flare với đôi mắt mở to, như thể vừa nghe một lời tiên tri xấu: ``Vậy là song trùng rồi còn gì! Chuẩn mô-típ kinh dị quá luôn, peko!'' cô lùi lại một bước, ôm chặt chiếc gối hơn như thể nó là vật phòng thân. ``Tôi sẽ không để ai thay thế tôi đâu, peko!''

Nhìn thấy tình hình có vẻ căng thẳng, Roboco lục trong túi, lấy ra một thiết bị nhỏ cỡ nắm tay, với chiếc đèn trắng nhấp nháy chậm rãi. Cô giơ nó lên và hất cằm về phía nhóm bạn, nở một nụ cười đầy bí ẩn: ``Đây là `Cảm biến ma' mà chị mua được ở chợ Akiba ảo hôm nọ này\~{}''

\img{5-11b}

Flare tò mò nghiêng đầu, đôi mắt sáng lên vì hứng thú: ``Nó để làm gì vậy? Trông như một món đồ chơi công nghệ cao ấy.''

Nàng Rô-bốt giải thích với giọng điệu đầy bí ẩn, hạ thấp giọng như thể đang tiết lộ một bí mật quan trọng: ``Nó sẽ nhấp nháy khi có ma ở gần. Càng nhiều ma, nó càng nhấp nháy nhanh. Đây là công nghệ mới nhất mà chị tìm hiểu được đó\~{}''

Pekora nheo mắt nhìn thiết bị một lúc lâu, rồi bĩu môi, giọng đầy vẻ hoài nghi: ``Nhìn như một món đồ chơi bán ở hội chợ vậy, peko. Em cá là nó chỉ là một cái đèn pin bình thường được gắn thêm vài cái đèn LED thôi.''

Roboco phì cười nhưng không đáp lại. Trong lúc đó, nàng Yêu tinh vàng quay sang nhìn A-chan, người vẫn đang chăm chú vào màn hình máy tính, hoàn toàn không để ý đến cuộc trò chuyện xung quanh.

Pekora tò mò hỏi, giọng lớn hơn một chút để cô nàng Đeo kính có thể nghe thấy: ``A-chan, chị không xem được phim kinh dị, đúng không?''

\dots Im lặng.

Không có hồi âm. Âm thanh duy nhất trong phòng lúc này là tiếng gõ bàn phím đều đặn. A-chan vẫn cúi đầu, mắt dán vào màn hình, ngón tay vẫn tiếp tục gõ những dòng lệnh. Thỉnh thoảng cô khựng lại một chút vì tưởng có người nào đó tới, nhưng rồi tiếp tục với công việc của mình, như thể cô đang ở trong một thế giới riêng.

Roboco nhún vai, kết luận đơn giản, giọng đầy vẻ hiểu biết: ``Vậy chị ấy cũng xem được nhỉ. Hoặc là chị ấy đã quá tập trung đến mức không nghe thấy gì. Cả hai đều có khả năng.''

Nhắc đến phim, Pekora vội lắc đầu nguầy nguậy. Cô cố gắng thuyết phục hai người họ, giọng như một người đang đưa ra một đề xuất rất hợp lý: ``Với một người trong sáng như Pekora thì thà xem phim tài liệu về chim cánh cụt ở Nam Cực còn hơn. Ít nhất thì chim cánh cụt không làm em sợ, peko.''

Chẳng biết từ lúc nào, Flare đã lục trong túi, rút ra một củ cà rốt sáng bóng, tươi rói. Cô đung đưa nó trước mặt Pekora, khẽ cười đầy ẩn ý, như một pháp sư đang chuẩn bị tung chiêu.

\img{5-11c}

Nàng Thỏ thoáng nhìn thấy màu cam thân thuộc thì mắt lập tức sáng lên. Đôi mắt cô dán chặt vào củ cà rốt, như thể bị thôi miên. Như bị một lực hút vô hình, cô chớp mắt vài lần, rồi bất giác buột miệng, giọng như một người đang nói trong mơ: ``Tự dưng tôi cảm thấy muốn xem phim kinh dị quá, peko.''

Flare nở nụ cười đầy thỏa mãn khi mọi thứ diễn ra đúng kế hoạch. Nàng Yêu tinh vàng vừa thành công thực hiện một cú ``điều khiển tâm trí'' đơn giản nhưng hiệu quả, và cô biết rằng từ bây giờ, Pekora sẽ không thể cưỡng lại sức hút của củ cà rốt, và cả bộ phim kinh dị sắp tới.

Roboco nhìn cảnh tượng, lắc đầu, nhưng cũng không giấu được nụ cười: ``Em không biết nên gọi đấy là ma thuật hay là thao túng tâm lý nữa.''

Flare cười khúc khích, vẫn đung đưa củ cà rốt: ``Gọi là gì cũng được, miễn là nó hiệu quả.''

\ct

Sau một hồi sắp xếp chỗ ngồi, Pekora, Roboco và Flare bật ti-vi lên và bắt đầu bộ phim của họ. Ánh sáng từ màn hình nhấp nháy trong căn phòng tối, hắt lên những cái bóng dài ngoằng trên tường, chúng uốn éo và nhảy múa theo từng khung hình như những linh hồn đang tìm cách thoát ra. Trong khi nàng Yêu tinh vàng và nàng Người máy có vẻ khá hào hứng với cốt truyện, nàng Thỏ vẫn co rúm lại, ôm chiếc gối nhỏ như một tấm bùa hộ mệnh. Cô thỉnh thoảng rụt cổ lại mỗi khi nhạc nền trở nên rùng rợn, đôi mắt cam mở to, sẵn sàng nhắm lại bất cứ lúc nào.

Nghe có động tĩnh phía sau, A-chan khẽ ngẩng đầu lên khỏi màn hình máy tính, liếc nhìn nhóm người sau lưng một lúc. Cô thấy Pekora đang ôm gối, Roboco đang vừa xem phim vừa nhai kẹo, và Flare đang chỉ tay vào màn hình dường như đang giải thích gì đó. Rồi cô lại cúi xuống tiếp tục công việc, không khỏi thầm nghĩ: ``\textit{Sao họ lại xem phim kinh dị vào lúc nửa đêm thế này?}''

``Để chị đi lấy chút đồ ăn vặt nha,'' Roboco chợt nói, đứng dậy khỏi ghế. Cô vươn vai, đôi mắt LED sáng lên trong bóng tối.

\img{5-11d}

Pekora lập tức bật dậy theo, vẻ mặt lộ rõ sự hoang mang. Cô túm lấy tay áo của Roboco, giọng như một đứa trẻ bị bỏ rơi: ``Chờ đã, em đi cùng nữa, peko!'' cô nói, như thể sợ ở lại đây một mình sẽ bị ám bởi bóng ma nào đó.

Nàng Người máy chớp mắt, chưa kịp phản ứng thì nàng Thỏ bỗng cảm thấy có một luồng khí lạnh thoảng qua từ phía cửa vào. Một cảm giác bất an len lỏi trong lòng, khiến tai cô dựng đứng lên và lông trên tay cũng dựng ngược.

\img{5-11e}

Pekora lắp bắp, chỉ tay về phía cửa, giọng như sắp khóc: ``C-Có cái gì ghê rợn (eerie) ở kia kìa!''

Roboco nghiêng đầu, chưa kịp hiểu. Cô nhìn về phía cánh cửa tối om, nhưng chỉ thấy bóng tối và những bóng cây bên ngoài cửa sổ: ``Eevee á? Em không thấy gì cả.''

Nàng Thỏ quát lên, giọng đầy bực bội: ``\textbf{Không! Không phải thứ mà chị bắt bằng cái Pékoball đâu! Là thứ gì đó khác!}

Hai người nhìn chằm chằm về phía hành lang tối om. Ở đó, có hai bóng người đang dần tiến về phía họ. Bước chân chậm rãi, hầu như không phát ra âm thanh, bóng dáng quen thuộc nhưng lại có gì đó\dots\ kỳ lạ. Cách họ di chuyển không giống người bình thường; nó uyển chuyển nhưng cũng có chút cứng nhắc, như một bộ phim chiếu chậm.

Pekora cứng đờ, hai tay bám chặt lấy cánh tay Roboco, giọng run rẩy: ``Ai đó đang ở kia kìa\dots''

Bóng người dần đi vào vùng sáng của căn phòng, và khuôn mặt họ lộ diện. Đó không ai khác chính là nàng thần tượng Sora và cậu Tổng quản lý Kuon. Gương mặt họ có vẻ bình thường, nhưng ánh mắt họ có một sắc thái kỳ lạ: không phải ngạc nhiên, không phải tò mò, mà là một vẻ trống rỗng đáng lo ngại, nhưng chỉ kéo dài trong tích tắc.

\img{5-11f2}

Sora nghiêng đầu, ánh mắt đầy sững sờ, nhưng có vẻ như cô đang cố tỏ ra tự nhiên: ``Ủa? Mọi người ở đây làm gì vậy? Chị tưởng giờ này chỉ có A-chan ở lại thôi chứ?''

Kuon khoanh tay, liếc nhìn nhóm người trong phòng rồi nhướng mày, nhưng giọng cậu có vẻ đều đều, thiếu đi sự tự nhiên thường thấy: ``Đêm hôm rồi mà mấy đứa chưa về à? Khuya lơ khuya lắc rồi đấy\dots''

Pekora và Roboco đồng loạt thở phào nhẹ nhõm, cảm giác sợ hãi ban nãy như tan biến đi một nửa. Nhưng sự kỳ lạ vẫn còn đọng lại trong không khí: một chút gì đó không được ổn.

Nàng Người máy lên tiếng trước, vẫn còn chút ngờ vực: ``Sao hai người lại ở đây? Em tưởng tối nay mọi người đều đã về hết rồi.''

Nàng Thỏ chớp mắt, vẫn còn chưa hoàn toàn định thần, tay vẫn nắm chặt áo Roboco: ``Sora-senpai? Kuon-senpai? Sao hai anh chị tới đây vậy?''

Nàng Thần tượng cười hiền, nhưng nụ cười của cô có vẻ hơi cứng nhắc: ``Chị đang chờ A-chan làm xong việc. Chị có hẹn với cậu ấy từ chiều, nhưng cậu ấy bảo có việc gấp nên chị đành chờ ở đây.''

Kuon nhún vai, giọng lười biếng như thường lệ, nhưng vẫn thiếu đi sự thoải mái thông thường: ``Anh thì bị Sora-chan rủ đi. Tiện thể xem tiến độ công việc của A-san. Anh cũng có vài thứ cần hỏi chị ấy đây.''

Những lời nói quen thuộc ấy lẽ ra phải khiến mọi người yên tâm hơn, nhưng không hiểu sao, bầu không khí trong phòng vẫn có chút kỳ lạ. Có một sự căng thẳng vô hình đang len lỏi trong căn phòng, như thể không khí đã trở nên đặc quánh hơn một chút.

Dường như ai đó trong nhóm vẫn cảm thấy có gì đó\dots\ không ổn. Flare, người đã quan sát họ từ đầu, cảm thấy một sự ngờ vực nhỏ nhoi đang lớn dần trong lòng.

Hai người vừa mới đến đứng yên lặng, không nói gì thêm. Sự im lặng kéo dài hơn một chút so với bình thường, và nó trở nên nặng nề đến mức Flare cuối cùng đã phải lên tiếng.

Nàng Yêu tinh hơi nheo mắt, quan sát họ thật kỹ cách họ đứng, cách họ cầm tay, cách họ nhìn vào mắt người khác. Cuối cùng, cô lên tiếng, cố giữ giọng bình thản: ``Có chuyện gì vậy? Hai người trông có vẻ hơi khác thường đấy.''

Kuon ôm trán, thở dài một cách quen thuộc: ``Cái câu ấy phải là bọn anh hỏi mới đúng đó. A-san đâu? Chị ấy vẫn ở đây chứ?''

Nghe rằng có ai đó nhắc đến mình, A-chan cuối cùng cũng chịu ngẩng đầu lên. Cô quay về phía hai người mới đến, đẩy gọng kính một cách thản nhiên, giọng vẫn đều đều: ``Mấy đứa gặp được ai à? Hay là gặp ma ở ngoài hành lang?''

Pekora vội vàng lắc đầu, cố gắng trấn tĩnh lại, nhưng tay vẫn không rời khỏi Roboco: ``Bọn em ngạc nhiên vì Sora-senpai và Kuon-senpai ở đây thôi ạ, peko. Tưởng tối nay chỉ có mình bọn em với A-chan thôi.''

Sora chớp mắt, rồi bật cười: ``À, chị cũng ngạc nhiên vì gặp mọi người ở đây mà. Chị tưởng tối nay mọi người sẽ về nhà sớm.''

Kuon khoanh tay, dựa người vào tường, tạo dáng như một người đang thư giãn: ``Nhất là vào giờ này. Mọi người nên về nghỉ ngơi đi, đừng có thức khuya mãi.'' Đoạn, cậu hướng ánh nhìn về phía bàn làm việc của A-chan ``Mà A-san, cái báo cáo em gửi hôm trước ấy, chị làm đến đâu rồi ạ? Có tiến triển gì không?''

A-chan nhìn vào màn hình, tay vẫn gõ bàn phím đều đặn, như thể cô đang ở trong một thế giới riêng. Cô đáp bằng giọng dửng dưng, không ngẩng đầu lên: ``Đợi tôi một tí. Sắp xong rồi đây. Cứ biết như vậy. Còn khoảng mười dòng nữa.''

Tiếng gõ phím vang lên trong không gian. Nhưng dù mọi thứ có vẻ bình thường, cảm giác kỳ lạ vẫn chưa hoàn toàn tan biến. Nó vẫn lơ lửng trong không khí, như một bóng ma vô hình đang quan sát tất cả bọn họ.

Roboco nắm lấy tay của Sora, giọng đầy phấn khích, cố gắng phá tan bầu không khí căng thẳng: ``Bà đã đến đây rồi, sao không xem phim với bọn tôi nhỉ? Phim này hay lắm, mới ra mắt tuần trước đấy.''

Pekora cũng lập tức phụ họa, ánh mắt long lanh đầy hy vọng: ``Phải đó! Ở lại xem với bọn em đi, peko! Sẽ vui hơn nhiều nếu có Sora-senpai và Kuon-senpai nữa!''

Cả hai người mới đến nhìn nhau, có vẻ đang suy nghĩ. Trong khi đó, hai cô nàng vẫn tiếp tục nài nỉ, mỗi người một cách. Có gì đó thật sự không ổn ở đây, nhưng không ai có thể xác định được điều gì.

``Xin hai người đó, peko! Nếu xem cùng nhau thì sẽ không thấy sợ nữa!'' nàng Thỏ tiếp tục năn nỉ, giọng như một đứa trẻ đang cầu xin bố mẹ cho phép xem phim.

Nàng Rô-bốt cũng gật gù, đồng tình với quan điểm đó, đôi mắt LED sáng lên đầy hào hứng: ``Đúng đấy! Phim hay hơn khi xem chung, và bọn em đã có snack sẵn rồi.''

Nhưng dù cho những lời thuyết phục có tha thiết đến đâu, Sora vẫn lắc đầu, nở một nụ cười áy náy: ``Xin lỗi, chị không thể xem với mấy đứa được. Chị chỉ đến đây để đón A-chan thôi. Cậu ấy đã hẹn chị sẽ đưa về rồi.''

Kuon cũng khoanh tay, nhẹ nhàng từ chối, nhưng cậu không nhìn thẳng vào mắt họ: ``Anh cũng không ở lại đâu. Qua một tí xong rồi về ngay ấy mà. Còn đống việc ở nhà đang chờ.''

Nét mặt hai cô gái trước mặt lập tức xị xuống. Pekora thậm chí còn lộ rõ vẻ vừa thất vọng vừa sợ hãi, đôi tai thỏ cụp xuống như một chú thỏ bị bỏ rơi. Nhưng chưa kịp để nàng Thỏ tiếp tục than thở, một tiếng thở dài vang lên từ góc phòng.

\img{5-11g}

A-chan lảo đảo bước đến, đôi mắt lờ đờ vì thiếu ngủ. Cô dụi mắt, vươn vai, và nói với giọng của một người vừa hoàn thành một cuộc chạy bộ: ``Tôi xong rồi đây\dots\ Cuối cùng cũng xong. Ai mà nghĩ là một báo cáo có thể dài đến thế.''

Sora và Kuon đồng loạt lên tiếng, gần như theo phản xạ, như thể đã chuẩn bị sẵn câu nói này: ``Chị làm việc vất vả rồi.''

Nàng Đeo kính vươn vai, ngáp dài một cái rồi nhìn xung quanh, chợt nhận ra điều gì đó. Mắt cô dừng lại ở chiếc TV vẫn đang sáng, phát những khung hình đầy bóng tối và tiếng nhạc rùng rợn: ``À mà, có phải hai người bật TV lên không, Sora, Hikawa-san?''

Hai người nhìn cô, vẻ khó hiểu. Họ đồng thanh với giọng đầy ngạc nhiên: ``Không. Bọn em vừa mới đến, chưa kịp chạm vào điều khiển.''

Ba người quay đầu lại nhìn về phía sau. Trong căn phòng bỗng chốc trở nên im ắng lạ thường. Màn hình TV vẫn đang phát sáng, nhấp nháy những khung hình kinh dị, nhưng không có ai đứng gần đó cả. Không một bóng người. Nhưng chiếc TV vẫn sáng.

\img{5-11h}

Sora nhíu mày, giọng có chút lo lắng: ``Hay là chị giẫm vào cái điều khiển rồi?''

A-chan khoanh tay, vẻ mặt có chút không vui, giọng như một người vừa bị nghi ngờ: ``Nếu thế thì tôi đã phải biết rồi chứ! Tôi đang làm việc, lấy đâu ra thời gian để bật TV?''

Kuon nhìn quanh một lượt, rồi lơ đãng phỏng đoán, cố tìm một lời giải thích hợp lý: ``Cũng có thể là đồng hồ hẹn giờ chưa tắt, chắc thế. Có ai đặt hẹn giờ bật TV không?''

A-chan nhíu mày, lắc đầu: ``Tôi không nhớ có ai đặt hẹn giờ đâu. Họ còn chẳng biết chỉnh cài đặt ở chỗ nào nữa là.''

Cậu con trai nhặt bộ điều khiển trên bàn lên, bấm nút tắt TV. Màn hình tối sầm lại, căn phòng chỉ còn ánh sáng nhàn nhạt từ ánh cửa sổ và màn hình máy tính. Một sự im lặng bao trùm căn phòng, im lặng đến mức có thể nghe thấy tiếng tim đập của mỗi người.

Nàng Đeo kính nhìn chằm chằm vào chiếc TV tối màu, như thể đang chờ nó tự bật lại. Cô lẩm bẩm, giọng đầy suy tư: ``Có lẽ chỉ là trục trặc kỹ thuật thôi\dots''

\ct

Sora bỗng bật cười khẽ, giọng nhẹ như gió thoảng qua những hàng cây ngoài cửa sổ, phá tan bầu không khí im lặng vừa rồi: ``Mùa hè đúng là vui thật nhỉ. Những đêm dài, những câu chuyện, và những tiếng cười vang dưới ánh trăng.''

Kuon cũng cười, nhưng trong giọng nói lại có chút gì đó ẩn ý - một sự hiểu biết mà chỉ những ai đã trải qua đủ nhiều đêm ở công ty này mới có: ``Nếu nói về những câu chuyện ma quỷ thì đúng là như vậy. Ở đây, chẳng cần phải đến tháng Bảy mới có chuyện ma.''

A-chan lập tức rùng mình, ôm hai tay trước ngực như thể tự trấn an. Cô nhìn Kuon với ánh mắt đầy bất mãn: ``Kinh thế! Sao cậu lại nói vậy? Nghe cậu kể chuyện ma như thế làm tôi sợ hết hồn.''

Ba người tiếp tục cuộc trò chuyện, giọng nói dần trở nên vui vẻ hơn, như thể bóng ma của chiếc TV tự bật đã bị lãng quên. Tiếng cười của Sora vang lên trong căn phòng, hòa cùng tiếng gõ phím lách cách của A-chan - cô đã quay lại bàn làm việc, nhưng vẫn tham gia vào cuộc trò chuyện.

Chẳng ai để ý rằng ngay góc phòng, ba bóng người mờ nhạt vẫn đang đứng đó từ nãy đến giờ.

Những kẻ song trùng - bản sao của Pekora, Roboco và Flare - im lặng quan sát, ánh mắt vô hồn như những hình nhân rỗng tuếch. Chúng đứng bất động, không thở, không chớp mắt, chỉ nhìn chằm chằm vào ba người đang trò chuyện với vẻ mặt tĩnh lặng đến rợn người.

\img{5-11i}

Bản sao của nàng Người máy cuối cùng cũng lên tiếng, giọng nói đơn điệu như một chiếc máy móc không cảm xúc, không có chút nhấn nháy nào của con người: ``\emph{Tiếc ghê, ta không bắt được họ rồi.}''

Nó nhìn những người đang nói chuyện, đôi mắt LED của bản sao lóe lên một tia sáng xanh lam rồi tắt hẳn. Hai bản sao còn lại cũng chỉ khẽ gật đầu, như thể đã chấp nhận sự thất bại.

Cả ba bóng người dần nhạt đi, như những vệt khói bị gió cuốn, rồi biến mất hoàn toàn, như thể chưa từng tồn tại. Không một tiếng động, không một dấu vết, chỉ còn lại một khoảng không trống rỗng nơi chúng vừa đứng.

Bên dưới sàn, nơi mà Roboco đã để lại thiết bị cảm biến ma, một vòng đèn LED màu đỏ nhấp nháy dữ dội, lêu lên những tiếng bíp bíp liên hồi, như thể đang báo động về một mối nguy hiểm đã qua. Rồi chợt, tiếng bíp lịm hẳn, đèn trắng tắt ngấm, và thiết bị trở nên im lìm như một món đồ chơi vô tri.

\img{5-11j}

Ngoài cửa sổ, làn gió đêm khẽ thổi, mang theo một tiếng thì thầm mơ hồ. Có thể là từ những tán cây, có thể là từ một nơi nào đó xa hơn. Nhưng trong phòng, không ai nghe thấy gì ngoài tiếng cười của họ.

Bản sao đã biến mất. Nhưng câu hỏi vẫn còn đó: liệu chúng sẽ quay lại vào một đêm khác, khi không ai để ý?

%==========================================TRUYỆN PHỤ=========================================
\omk

``Và\dots\ cắt!''

Đèn trong căn phòng bật sáng trưng, xua tan bóng tối và những bóng ma vô hình. Từ đằng sau một trong những khu máy tính, Shion ló mặt ra với vẻ đầy tự hào, trên môi nở nụ cười đắc thắng. Cả nhóm tụ tập lại để xem thành quả của mình, những gương mặt rạng rỡ vì háo hức.

Trên màn hình, đoạn phim cuối cùng đã được ghi hình hoàn chỉnh. Hiệu ứng ánh sáng mờ ảo, tiếng nhạc nền rờn rợn, cộng với biểu cảm tự nhiên của dàn diễn viên, tất cả đều hòa quyện tạo nên một khung cảnh đậm chất kinh dị. Những bóng đen uốn éo trên tường, ánh mắt vô hồn của các bản sao, tất cả đều được tái hiện một cách hoàn hảo.

Yamazu và A-chan chăm chú theo dõi từng, khung hình, thỉnh thoảng trao đổi với nhau vài câu về góc quay và ánh sáng. Tay cậu bạn Hội họa lia lịa trên bàn phím, tua đi tua lại một vài phân đoạn để kiểm tra chi tiết. Sau một lúc trầm ngâm, hai người cùng gật đầu.

A-chan đẩy gọng kính lên, giọng đầy hài lòng: ``Được rồi. Cảnh này đạt yêu cầu. Ánh sáng và bóng đổ đều hoàn hảo.''

Yamazu gật gù, vẫn với giọng vấp quen thuộc nhưng không kém phần hào hứng: ``U-Ừ\dots\ Thế này là ổn rồi đó! C-camera đã bắt được hết các góc cần thiết.''

Vừa nghe thấy thế, Shion, Pekora, Roboco và Flare đồng loạt nhảy cẫng lên vui sướng. Họ ôm nhau, vỗ tay, và reo hò như những đứa trẻ vừa hoàn thành một dự án lớn.

Nàng Thỏ reo lên, đôi tai thỏ vểnh cao vì phấn khích: ``Thành công rồi! Bọn em làm được rồi, peko! Cuối cùng cũng xong!''

Nàng Người máy cũng không giấu nổi sự vui mừng, đôi mắt LED sáng rực như đèn pha: ``Chị biết ngay mà, chị diễn hoàn hảo lắm mà! Không uổng công em đã tập luyện cả tuần.''

Nữ Yêu tinh khẽ rùng mình sau khi xem lại cảnh quay của chính mình, tay ôm lấy cánh tay: ``Không ngờ là lại đáng sợ như thế luôn\dots\ Kế hoạch của Shion-senpai thành công rồi đấy! Cái ánh mắt của bản sao trông thật sự rợn người.''

Shion khoanh tay, hất cằm đầy kiêu hãnh, giọng như một bậc thầy phép thuật đang nhận lời khen: ``Tất nhiên rồi! Đây chính là ma thuật của phù thủy vĩ đại Shion-sama đó nha! Không có gì là không thể với chị cả!''

Trong khi đó, Pekora và Roboco ôm nhau như thể hai người vừa trải qua một trận chiến căng thẳng. Flare cũng mỉm cười nhẹ nhõm, rõ ràng là cô đã có chút lo lắng về việc bộ phim có đủ đáng sợ hay không, nhưng kết quả cuối cùng lại vượt xa mong đợi.

Kuon đứng gần đó, hai tay khoanh trước ngực, mỉm cười khi thấy mọi người đều hào hứng như vậy. Cậu quay sang Yamazu và hỏi: ``\vie{Tổng kết lại thì sao rồi? Có cần quay lại cảnh nào không?}''

Yamazu nói với những lời nói vấp quen thuộc, nhưng ánh mắt thì sáng lên vì tự hào: ``V-Về cơ bản\dots\ l-là mọi người diễn rất tốt rồi, c-camera thì đã quay từ nhiều góc khác nhau, n-nên ta c-cứ chọn góc tốt nhất rồi g-ghép vào là xong thôi! Không cần quay lại gì cả.''

``Vậy là phần hậu kỳ phải nhờ hai người rồi,'' Kuon nhìn cậu bạn và A-chan, cười trừ. ``Các cậu có thể xử lý được không?''

Cả Yamazu và A-chan cùng gật đầu, giọng đầy tự tin: ``Không thành vấn đề.''

Không khí trong phòng vẫn tràn ngập sự phấn khích khi cả nhóm dần bình tĩnh lại. Ở gần đó, dàn diễn viên của bộ phim vừa quay xong vẫn đang trò chuyện rôm rả.

Roboco vỗ vai Shion, giọng đầy ngưỡng mộ: ``Em giỏi thật đó, Shion-chan! Làm ra được cái thứ này không dễ đâu nha! Bọn chị cứ tưởng là ma thật đấy.''

Nàng Phù thủy ưỡn ngực tự hào, giọng như một nghệ sĩ vừa có buổi biểu diễn thành công: ``He he! Thế nào? Em chỉ vừa học phép này cách đây ít hôm đấy! Ma thuật của em luôn luôn tiến bộ từng ngày.''

Flare mỉm cười, gật gù: ``Chị làm lại cho bọn em xem đi. Bọn em muốn thấy cách tạo ra bản sao một lần nữa.''

Shion gật đầu đầy nhiệt tình: ``Được thôi!''

Nàng Phù thủy tóc xám vung tay lên, niệm một câu thần chú nhỏ. Ngay lập tức, ba bản sao của Pekora, Roboco và Flare hiện ra ngay trước mắt họ, đứng im như tượng.

Nàng Thỏ há hốc mồm, chớp mắt liên tục rồi lùi về sau một bước. Cô nhìn bản sao của mình với ánh mắt vừa kinh ngạc vừa sợ hãi: ``N-Nhìn y như đúc chúng ta luôn này, peko! Tóc, mắt, cả cách đứng nữa!''

Nàng Yêu tinh vàng nghiêng đầu nhìn bản sao của mình, ánh mắt đầy tò mò: ``Nhưng mà\dots\ sao trông chúng đờ đẫn thế? Mắt thì vô hồn, cử chỉ thì máy móc.''

``À thì\dots\ chị còn phải giữ sức cho bản thân nữa chứ! Để chị thêm chút ma lực vào là khác ngay.''

Cô vung tay niệm phép, và ngay lập tức ba bản sao có sức sống hẳn lên: ánh mắt có hồn hơn, dáng điệu tự nhiên hơn, thậm chí còn bắt chước được cử chỉ của bản gốc. Chúng bắt đầu di chuyển, nhìn quanh, và nhếch mép cười.

Pekora song trùng nghiêng đầu, nheo mắt nhìn bản gốc rồi bất ngờ lên tiếng với giọng điệu giống hệt cô, đầy thách thức: ``Sao bà nhìn giống tôi thế, peko? Hay là bà đang cố bắt chước tôi?''

Pekora tròn mắt, chỉ thẳng vào bản sao trước mặt mình, giọng đầy bực bội: ``Thì tôi là bà, peko! Tôi mới là người thật!''

Bản sao song trùng cười điệu ``Thỏ'', tay chống nạnh đầy tự tin, giọng y hệt: ``Không, tôi mới là bản xịn! Bà chỉ là bản sao thôi, peko!''

Pekora nhảy dựng lên, quơ quào tay loạn xạ như một con thỏ bị chọc tức: ``Không, là tôi mới đúng! Tôi có cà rốt, bà có không? Không, đúng không?''

Hai nàng Thỏ lao vào tranh cãi xem ai mới là Pekora thật, khiến mọi người xung quanh cười bò ra đất. Roboco thậm chí phải vịn vào vai Flare để khỏi ngã vì cười quá nhiều, trong khi bản thân Flare chỉ biết bật cười bất lực, nước mắt chảy ra vì vui.

Cuối cùng, Shion búng tay một cái, lập tức Pekora ``pha-ke'' trở lại trạng thái trước, đứng im lặng như một con búp bê vô hồn, đôi mắt trống rỗng nhìn về phía trước.

Sau đó, cô nàng phù thủy lại niệm phép một lần nữa, ba bản sao dần dần tan biến trong không khí như những làn khói mỏng, để lại một khoảng trống trước mắt cả nhóm. Chúng biến mất không một tiếng động, chỉ còn lại một cảm giác kỳ lạ như thể chúng chưa từng tồn tại.

Shion vừa phủi tay vừa cười đắc ý, giọng như một nhà ảo thuật vừa hoàn thành tiết mục: ``Thế nào? Quá chuẩn bài luôn chứ gì? Ma thuật của em không có gì phải bàn cãi.''

Flare khẽ cười, thỉnh thoảng run lên bần bật vì cười, giọng vẫn còn vui: ``Ừm\dots\ công nhận là nhìn ghê rợn thật đấy. Nếu không biết là giả, em đã tưởng là ma thật rồi.''

Roboco cười lớn, đưa tay lau nước mắt: ``Vậy mới gọi là phim kinh dị mùa hè chứ! Có chút ma quái mới đúng chất!''

Cùng lúc đó, Kuon và Sora vừa bước vào phòng, trên tay Sora cầm theo một khay nước mát cho mọi người. Những cốc nước đá lấp lánh dưới ánh đèn, mang theo hơi mát xua tan cái nóng của mùa hè.

``Mọi người đã cố gắng rồi!'' Sora vui vẻ nói, đặt khay nước xuống bàn. ``Mọi người uống nước đi, bù lại sức đã mất.''

Theo sau là Kuon vỗ tay, giọng đầy khích lệ: ``Chúc mừng mọi người! Bộ phim đã hoàn thành tốt đẹp. Một tác phẩm đáng tự hào.''

Cả năm thành viên vừa quay xong bộ phim cúi đầu cảm ơn, gương mặt ai nấy đều rạng rỡ. Kuon phì cười, xua tay ra hiệu rằng họ không cần khách sáo như vậy: ``Không, anh nghĩ là mấy đứa nên cảm ơn Coco, Towa và Okayu mới đúng á. Chính ba người đó cùng với Hokuro đã viết kịch bản cho số này.''

Cả nhóm ngớ mặt ra, chưa kịp hiểu ý câu nói của cậu. Shion chớp mắt vài lần, Pekora nghiêng đầu ngạc nhiên, còn Flare thì há hốc miệng. Kuon tiếp tục giải thích, giọng như một người tiết lộ một bí mật thú vị: ``Chính ba người đó cùng với Hokuro-san viết kịch bản cho số này mà. Coco đảm nhận phần hài, Towa phần kinh dị, còn Okayu và Hokuro lo phần kỹ thuật.''

Sora gật đầu, cười nhẹ, giọng ấm áp: ``Và sau thành công từ tập kinh dị năm ngoái, lần này họ cho chị vào đóng đó! Một vai khách mời nhỏ, nhưng rất vui.''

Bốn cô gái đồng loạt ``Ồ\dots'' rồi ``À\dots'' một tiếng, ánh mắt không giấu được sự bất ngờ. Họ nhìn nhau, rồi lại nhìn Kuon, như thể vừa nhận ra rằng đằng sau mỗi bộ phim đều có một đội ngũ hậu trường đầy tâm huyết.

``Mà, anh nghĩ mấy đứa cũng mệt rồi.'' cânh hướng về chiếc vô tuyến mới mua kia, mắt sáng lên một tia tinh nghịch. ``Hay là\dots''

Sora mở lời đề nghị tới mọi người, giọng như một đứa trẻ đang rủ bạn bè chơi trò mới: ``Ta cùng xem phim nha? Hôm nay là đêm chiếu phim marathon mà.''

Cả nhóm đồng loạt hưởng ứng, dù rằng vẫn chưa ai biết bộ phim đó sẽ là gì. Những tiếng ``Đồng ý!'', ``Hay quá!'', ``Chọn phim gì đi!'' vang lên khắp phòng.

Flare tò mò hỏi, nghiêng đầu nhìn về phía TV: ``Mọi người định xem phim gì vậy? Có ai có đề xuất không?''

Pekora thì cụp tai xuống, khẽ run run, như một chú thỏ nhỏ đang sợ hãi: ``Làm ơn đừng là phim kinh dị, peko\dots\ Em không xem nổi nữa đâu.''

Kuon nở một nụ cười nửa miệng, đáp lại bằng giọng đầy ẩn ý, như thể đang trêu chọc nàng Thỏ: ``Có gì đâu, chỉ là một bộ phim về một con người có hai mươi tư nhân cách khác nhau thôi mà.''

Cô nàng bị chọc lập tức biến sắc, hai tai giật giật như một cái radar báo động. Cô nhìn Kuon với vẻ mặt đầy khiếp sợ: ``H-Hả\dots!? Thế thì nó là phim kinh dị rồi còn gì!? Em không xem đâu, peko!''

Cả nhóm cười phá lên, tiếng cười vang vọng khắp căn phòng. Shion vỗ vai Pekora, Flare cười rung vai, còn Roboco thì thậm chí còn nhảy nhót theo điệu cười. Pekora thì phồng má, khoanh tay tỏ vẻ khó chịu, nhưng khóe miệng lại không giấu được nụ cười: ``Không vui tí nào đâu\dots\ peko\dots''

Không ai bận tâm đến lời phản đối yếu ớt đó, và một buổi ma-ra-tông phim bắt đầu. Mỗi người đều chọn một bộ phim mình thích, và họ thay phiên nhau xem liên tục cho đến tận đêm khuya.

Ánh sáng từ màn hình hắt lên những khuôn mặt mệt mỏi nhưng đầy thích thú. Trên bàn, đống đồ ăn vặt chất thành đống: từ bỏng ngô, kẹo, bim bim, tới cả những thanh sô-cô-la mà Roboco đã giấu trong người. Những lon nước giải khát trống rỗng nằm ngổn ngang, và Sora khẽ ngáp khi đồng hồ đã điểm gần nửa đêm.

Bên ngoài, một cơn gió mùa hè khẽ lùa qua cửa sổ, làm lay động những tấm rèm. Dù bộ phim chỉ là giả tưởng, nhưng không khí ma quái mà nó mang lại vẫn khiến ai đó khẽ rùng mình.

Ở góc phòng, Kuon khoanh tay tựa lưng vào ghế, chậm rãi cười. Ánh mắt cậu lướt qua từng khuôn mặt đang say sưa với bộ phim – những người bạn, những đồng nghiệp, những thành viên của một gia đình kỳ lạ mà cậu đã học cách yêu thương.

``Mùa hè đúng là vui thật nhỉ?''

Giống như một câu hỏi tu từ, hay cũng có thể là một lời nhắn nhủ tinh nghịch gửi đến bất kỳ ai đang đọc câu chuyện này. Bên ngoài cửa sổ, những vì sao vẫn đang lấp lánh, và trong căn phòng, tiếng cười vẫn tiếp tục vang lên.

\end{document}